% successor-v10 finite levels review, 9 October 2026.
% Candidate replacement/extension at Proposition 6.32 and Definition 6.29;
% not applied to the author's v10 main.tex. Merge with existing TODOs
% and preserve the four mathematical v10 repairs without editorial changes.
%
\todo[inline]{\textbf{v10 validation (finite levels of the partial-type tree).}
A complete $L^+$-type over $d$ socle vertices has only
finitely many values in the finite unary/binary signature:
it is determined by all directed binary, unary and diagonal
$L$-bits, together with its full auxiliary $E$-matrix.
Thus the raw tree of complete induced partial types has finite
levels, and the admissible family of forbidden-free
partial types inherits finite levels. The corresponding
statements are formalised in
\texttt{SuccessorTree/V10/FiniteKptLevels.lean}
(pending the exact-head Lean audit), without assuming
finiteness as a new axiom or treating arbitrary raw
types as admissible. To finish the literal (finite-level)
$K^\mathcal F_{pt}$ tree representation and its
\texttt{LevelTree} instance, combine this result with
the checked prefix closure and prove that the greatest
common raw prefix of any two admissible types with a
common root is itself admissible. The concrete
successor-operation axioms and $M$-monoid then remain
a separate application obligation.}
%
% Expected Lean endpoints:
%   rawPartialTypeLevel_finite
%   admissibleRawTypeLevel_finite
